\subsection{切向复合律}

设 X 和 Y 是类为 \(C^{r}\) 的流形。我们知道（《Differentiable and Analytic Manifolds》R, 8.1.4），\(X \times Y\) 是类为 \(C^{r}\) 的流形，并且映射 \((\mathrm{T}(\mathrm{pr}_{1}), \mathrm{T}(\mathrm{pr}_{2}))\)，即到典范投影的切映射之积，是类为 \(C^{r-1}\) 的同构，从 \(\mathrm{T}(\mathrm{X} \times \mathrm{Y})\) 到 \(\mathrm{T}(\mathrm{X}) \times \mathrm{T}(\mathrm{Y})\)。这个同构与基空间为 \(X \times Y\) 的向量丛结构相容，并允许我们将 \(\mathrm{T}(\mathrm{X} \times \mathrm{Y})\) 识别为 \(\mathrm{T}(\mathrm{X}) \times \mathrm{T}(\mathrm{Y})\)。令 \(a \in \mathbf{X}, b \in \mathbf{Y}, u \in \mathrm{T}_a(\mathbf{X}), v \in \mathrm{T}_b(\mathbf{Y})\)；上述识别使我们可以把 \((u, v)\) 视为 \(\mathrm{T}_{(a,b)}(\mathbf{X} \times \mathbf{Y})\) 中的元素；于是

\[
(u, v) = (u, 0) + (0, v)
\]

并且 \((u,0)\)（相应地，\((0,v))\) 是 \(u\)（相应地，\(v\)）在浸入 \(x\mapsto (x,b)\)（相应地，\(y\mapsto (a,y)\)）的切映射下的像，该浸入将 X（相应地，Y）映入 \(\mathbf{X}\times \mathbf{Y}\)。需要精确说明时，我们将 \(0_{a}\) 记为 \(\mathrm{T}_a(X)\) 的零元素。

现在设 X、Y、Z 是类为 \(C^{r}\) 的流形，且 f 是类为 \(C^{r}\) 的从 \(X \times Y\) 到 Z 的映射。利用上述识别，切映射是类为 \(C^{r-1}\) 的从 \(T(X) \times T(Y)\) 到 \(T(Z)\) 的映射。对于 \(u \in T_{a}(X)\) 和 \(v \in T_{b}(Y)\)，

\begin{enumerate}
\def\labelenumi{(\arabic{enumi})}
\tightlist
\item
\end{enumerate}

\[
\mathrm{T} (f) (u, v) = \mathrm{T} (f) (u, 0 _ {b}) + \mathrm{T} (f) (0 _ {a}, v),
\]

\[
\mathrm{T} (f) \left(0 _ {a}, 0 _ {b}\right) = 0 _ {f (a, b)}.\tag{2}
\]

另一方面，映射 \(y \mapsto f(a, y)\) 是浸入 \(y \mapsto (a, y)\) 与 f 的复合；因此

\begin{enumerate}
\def\labelenumi{(\arabic{enumi})}
\setcounter{enumi}{2}
\item
  \(\mathrm{T}(f)(0,v)\) 是 v 在切映射 \(y\mapsto f(a,y)\) 下的像。同样，
\item
  \(\mathbf{T}(f)(u,0)\) 是 \(u\) 在切映射 \(x\mapsto f(x,b)\) 下的像。
\end{enumerate}

若映射 \(f\) 从 \(\mathbf{X} \times \mathbf{Y}\) 到 \(\mathbf{Z}\) 记作 \((x, y) \mapsto xy\)，则常用 \(uv\) 表示元素 \(T(f)(u, v)\)，其中 \(u \in T(\mathbf{X})\)、\(v \in T(\mathbf{Y})\)。

设 X 是类为 \(C^{r}\) 的流形，且有复合律 \(m: X \times X \to X\)，其类为 \(C^{r}\)。则 \(T(m)\) 是类为 \(C^{r-1}\) 的 \(T(X)\) 上的复合律。它称为与 m 相切的复合律。从 \(T(X)\) 投影到 X 的典范投影 p 与复合律 m 和 \(T(m)\) 相容；换言之，

\[
p \circ \mathrm{T} (m) = m \circ (p \times p).\tag{5}
\]

由 (2) 可得

\[
\mathrm{T} (m) \left(0 _ {x}, 0 _ {y}\right) = 0 _ {m (x, y)}\tag{6}
\]

对所有 \(x, y\) 属于 \(\mathbf{X}\)，换言之，零截面 \(x \mapsto 0_x\)（属于 \(\mathrm{T}(\mathbf{X})\)）与律 \(m\) 和 \(\mathrm{T}(m)\) 相容。

命题 1. 设 X 是类为 \(\mathbf{C}^r\) 的流形，且 \(m\) 是 X 上类为 \(\mathbf{C}^r\) 的复合律。若 \(m\) 是结合的（相应地，交换的），则 \(\mathrm{T}(m)\) 是结合的（相应地，交换的）。

若 \(m\) 是结合的，则 \(m \circ (m \times \mathrm{Id}_{\mathbf{x}}) = m \circ (\mathrm{Id}_{\mathbf{x}} \times m)\)，从而

\[
\mathrm{T} (m) \circ (\mathrm{T} (m) \times \mathrm{Id} _ {\mathrm{T} (\mathbf {X})}) = \mathrm{T} (m) \circ (\mathrm{Id} _ {\mathrm{T} (\mathbf {X})} \times \mathrm{T} (m))
\]

从而 \(\mathrm{T}(m)\) 是结合的。令 \(s\) 为映射 \((x,y)\mapsto (y,x)\)，它从 \(\mathbf{X}\times \mathbf{X}\) 到 \(\mathbf{X}\times \mathbf{X}\)。若 \(m\) 是交换的，则 \(m\circ s = m\)，从而

\[
\mathrm{T} (m) \circ \mathrm{T} (s) = \mathrm{T} (m).
\]

但是，\(\mathrm{T}(s)\) 是映射 \((u,v)\mapsto (v,u)\)，它从 \(\mathrm{T}(\mathbf{X})\times \mathrm{T}(\mathbf{X})\) 到 \(\mathrm{T}(\mathbf{X})\times \mathrm{T}(\mathbf{X})\)。因此 \(\mathrm{T}(m)\) 是交换的。

命题 2. 设 X 是类为 \(\mathbf{C}^r\) 的流形，\(m\) 是 X 上类为 \(\mathbf{C}^r\) 的复合律，\(e\) 是 \(m\) 的单位元。

\begin{enumerate}
\def\labelenumi{(\roman{enumi})}
\item
  向量 \(0_{e}\) 是 \(\mathbf{T}(m)\) 的单位元。
\item
  \(\mathrm{T}_{e}(\mathbf{X})\) 在 \(\mathrm{T}(m)\) 下稳定，并且 \(\mathrm{T}_{e}(\mathbf{X})\) 上由 \(\mathrm{T}(m)\) 诱导的复合律是 \(\mathrm{T}_{e}(\mathbf{X})\) 上的向量空间加法。
\item
  令 U 为 X 的开子集，\(\alpha\) 为从 U 到 X 的类为 \(C^{r}\) 的映射，使得对所有 \(x \in U\)，\(\alpha(x)\) 是 x 在 m 下的逆元。则对所有 \(u \in T(U)\)，\(T(\alpha)u\) 是 u 在 \(T(m)\) 下的逆元。
\end{enumerate}

性质 (3) 和 (4) 表明，\(\mathrm{T}(m)(0_e,u) = \mathrm{T}(m)(u,0_e) = u\) 对所有 \(u\in \mathrm{T}(\mathbf{X})\) 成立，从而得到 (i)。对于 \(u,v\)，它们属于 \(\mathrm{T}_e(\mathbf{X})\)，

\[
\mathrm{T} (m) (u, v) = \mathrm{T} (m) (u, 0 _ {e}) + \mathrm{T} (m) (0 _ {e}, v) = u + v,
\]

从而得到 (ii)。最后，关系 \(m(x, \alpha(x)) = m(\alpha(x), x) = e\) 对所有 \(x \in \mathbf{U}\) 成立，这意味着

\[
\mathrm{T} (m) (u, \mathrm{T} (\alpha) (u)) = \mathrm{T} (m) (\mathrm{T} (\alpha) u, u) = 0 _ {e}
\]

对所有 \(u \in \mathrm{T}(\mathrm{U})\) 成立，从而得到 (iii)。

命题 3. 设 \(\mathbf{X}_1, \mathbf{X}_2, \ldots, \mathbf{X}_p\)、Y 是类为 \(\mathbf{C}^r\) 的流形，i 是 \([1, p]\) 中的整数，\(m_t\)（相应地，n）是类为 \(\mathbf{C}^r\) 的复合律，作用于 \(\mathbf{X}_t\)（相应地，Y），且 u 是类为 \(\mathbf{C}^r\) 的从 \(\mathbf{X}_1 \times \mathbf{X}_2 \times \dots \times \mathbf{X}_p\) 到 Y 的映射。若 u 关于下标为 i 的变量满足分配律，则 T(u) 关于下标为 i 的变量也满足分配律。

证明类似于命题 1 的证明。

